\chapter{Para Mantenedores: o Grafo de Depend{\^e}ncias e a Auditoria do Acervo}
\label{ap:mantenedor}

\begin{dicaopencode}
\textbf{Voc{\^e} n{\~a}o precisa conhecer o funcionamento interno das skills para utiliz{\'a}-las}: este ap{\^e}ndice descreve como o acervo {\'e} mantido e auditado ({\'u}til para quem deseja criar ou modificar skills). Para o uso em sala de aula, basta digitar \texttt{/nome-da-skill} ou mencionar o slug no \emph{prompt} -- o restante {\'e} transparente para o usu{\'a}rio.
\end{dicaopencode}

As 62 skills do reposit{\'o}rio n{\~a}o operam isoladamente; constituem um \textbf{grafo direcionado de conhecimento pedag{\'o}gico}.

Por exemplo, a skill de constru{\c c}{\~a}o de avalia{\c c}{\~o}es (\texttt{ia-educacao-avaliacao}) aponta depend{\^e}ncias formais para \texttt{ia-educacao-bloom} (para alinhamento cognitivo), \texttt{ia-educacao-rubrica} (para crit{\'e}rios de corre{\c c}{\~a}o), \texttt{aias-consultant} (para governan{\c c}a do uso de IA) e \texttt{ia-educacao-planejamento-reverso} (para enquadramento curricular).

Para assegurar que o grafo n{\~a}o se degenere com o crescimento do acervo, o reposit{\'o}rio conta com uma rotina de auditoria automatizada no script \texttt{audit.sh}:
\begin{itemize}
  \item \textbf{Zero Skills {\'O}rf{\~a}s}: Toda skill do acervo deve ser referenciada como depend{\^e}ncia por ao menos uma outra skill, garantindo relev{\^a}ncia sist{\^e}mica;
  \item \textbf{Zero Depend{\^e}ncias Quebradas}: Qualquer refer{\^e}ncia cruzada em \texttt{\#\# Depend{\^e}ncias} deve corresponder a uma pasta real e v{\'a}lida dentro de \texttt{skills/};
  \item \textbf{Isolamento de Depend{\^e}ncias Externas}: As skills s{\~a}o autocontidas neste reposit{\'o}rio: nenhuma referencia skills externas ou pacotes n{\~a}o controlados;
  \item \textbf{Conformidade Terminol{\'o}gica}: O script verifica a consist{\^e}ncia exata dos nomes dos 5 n{\'i}veis da Escala AIAS e a grafia institucional em refer{\^e}ncias bibliogr{\'a}ficas normativas do MEC.
\end{itemize}
